% ═══════════════════════════════════════════════════════════════════════════
% Master-thesis chapter OUTLINE: Results
% ---------------------------------------------------------------------------
% Skeleton only — bullet points state WHAT to report and WHICH figure/table
% carries it; no result descriptions (numbers unknown until the full batch
% runs). Every figure placeholder names the notebook + export that produces
% it, and — where it already exists — the file in Figures_paper/.
%
% Placeholder macro: renders a framed box so the document compiles before the
% figures exist. Swap \figplaceholder{...} for \includegraphics once ready.
% ═══════════════════════════════════════════════════════════════════════════

\newcommand{\figplaceholder}[2][0.85\textwidth]{%
  \fbox{\parbox[c][0.35\textwidth][c]{#1}{\centering\ttfamily #2}}}

\section{Results}
\label{sec:results}

% Suggested running order: static → quasi-static → dynamic → amplitudes →
% linearity → asymmetry → modal. Mirrors the processing chapter, and builds
% the physical story: geometry sets Θ, Θ predicts the quasi-static transfer,
% the resonance structures the dynamic transfer, and the last three sections
% probe the model's assumptions (linearity, symmetry, mode identity).

% ─────────────────────────────────────────────────────────────────────────────
\subsection{Static Configuration: Geometry and Sag}
\label{sec:res:static}
% Source: notebook 01 → geometry_sag.npz

\begin{itemize}
  \item Table of all measured configurations: cable, gap length, sag variant,
    chord length, measured sag $w_0$ ($\pm$ parabola-fit residual),
    $\Theta_\mathrm{sag}$, $\Theta_\mathrm{mat}$ ($\pm$ std from the three
    $E$ tests), $\eta_\mathrm{pred} = 1/(1+\Theta)$.
    % → table from geometry_sag.npz; candidate for the appendix if too long,
    %   with a condensed per-cable version in the text.
  \item Measured vs.\ tension-free theoretical sag: which cables/gaps lie on
    the analytical line, which fall below (residual pretension from
    mounting), and whether the intended-sag variants reach their target sag.
  \item Consequence for the parameter space: the range of $\Theta$ actually
    spanned by the experiment (ideally several decades, straddling the
    $\Theta = 1/32$ rod-to-string transition).
  \item Comparison $\Theta_\mathrm{sag}$ vs.\ $\Theta_\mathrm{mat}$: since
    $\Theta_\mathrm{mat} \propto L^8$, gap-length uncertainty makes it far
    less reliable — justifies using $\Theta_\mathrm{sag}$ on all following
    x-axes (ties back to the uncertainty section of the processing chapter).
\end{itemize}

\begin{figure}[htb]
  \centering
  \figplaceholder{sag\_measured\_vs\_theory (exists in Figures\_paper/)\\
    measured $w_0$ vs.\ $\rho A g L^4 / 4\pi^4 EI$, log–log, 1:1 line,\\
    colour = cable, marker = gap, open = intended-sag variant}
  \caption{Measured midpoint sag against the tension-free analytical
    prediction. [Describe: proximity to 1:1 line, pretension offsets,
    effectiveness of the intended-sag variants.]}
  \label{fig:res:sag}
\end{figure}

% Optional second static figure (from plotting.plot_sag_vs_force_ratio):
% dimensionless sag w0/L vs force ratio, the Probst Fig. 2a reproduction —
% nice if you want to mirror the paper's figure sequence.

% ─────────────────────────────────────────────────────────────────────────────
\subsection{Quasi-Static Strain Transfer: $\eta$ vs.\ $\Theta$}
\label{sec:res:etatheta}
% Source: notebook 03 §6 → freqdomain_eta.npz. THE headline figure.

\begin{itemize}
  \item Headline result: coherence-gated band-mean $\eta$ (5–25 Hz) for all
    $\sim$42 configurations against $\Theta_\mathrm{sag}$, overlaid with the
    analytical $\eta = 1/(1+\Theta)$.
  \item Report: agreement/deviation from the model curve per $\Theta$ regime;
    behaviour around the $\Theta = 1/32$ threshold; whether taut
    configurations cluster near 100\,\%.
  \item Systematics to comment on: cable-to-cable trends at similar $\Theta$
    (material effects beyond the model?), gap-length trends, taut vs.\
    intended-sag pairs of the same cable (the controlled $\Theta$ increase).
  \item Error discussion: vertical bars (std over coherent band bins,
    H1/H2 bracket) and horizontal bars (sag residual → $\Theta$).
  \item Optional companion: same scatter with $\Theta_\mathrm{mat}$ on the
    x-axis, showing the larger horizontal scatter — visual argument for the
    sag-based parameterisation.
\end{itemize}

\begin{figure}[htb]
  \centering
  \figplaceholder{eta\_vs\_theta (exists in Figures\_paper/)\\
    $\eta$ [\%] vs.\ $\Theta_\mathrm{sag}$, log x, theory curve
    $100/(1+\Theta)$,\\ $\Theta = 1/32$ marker line}
  \caption{Quasi-static strain-transfer efficiency against the sag parameter
    $\Theta$ for all configurations. [Describe: model agreement, threshold
    behaviour, outliers.]}
  \label{fig:res:etatheta}
\end{figure}

% ─────────────────────────────────────────────────────────────────────────────
\subsection{Dynamic Strain Transfer: Frequency Response and Resonance}
\label{sec:res:dynamic}
% Source: notebook 03 §2–3 → freqdomain_frf.npz

\begin{itemize}
  \item Four-panel FRF for 2–4 representative configurations spanning a range
    of $\Theta$ (e.g.\ one taut, one sagged pair of the same cable):
    midpoint amplitude + phase, strain-transfer amplitude + phase.
  \item The expected regime structure (cf.\ theory chapter / Simone Fig.~3):
    flat quasi-static plateau at $1/(1+\Theta)$ → dip just below $f_1$
    (destructive coupling) → amplification spike above $f_1$ (constructive
    coupling, $>$100\,\%) → recovery towards 100\,\% at high frequency.
    Report which of these features the measurements show.
  \item The $\approx 180^\circ$ midpoint-phase flip across the resonance —
    the mechanism's fingerprint (wrapped-phase panels).
  \item Resonance table across configurations: $f_1$ measured, SDOF-fit $Q$
    and $\zeta$, theoretical clamped–clamped $f_1^\mathrm{th}$; discuss the
    systematic offset (rig/added-mass vs.\ ideal clamped beam) and the trend
    of $f_1$ with $\Theta$ (sag stiffening).
  \item Coherence: typical $\gamma^2$ per band; where the FRF is trustworthy
    and where (anti-resonances, low excitation) it is not.
\end{itemize}

\begin{figure}[htb]
  \centering
  \figplaceholder{frf\_four\_panel (exists in Figures\_paper/)\\
    2×2: $|H_\mathrm{mid}|$ norm., wrapped phase, $|H_\eta|$ [\%],
    wrapped phase;\\ $f_1$ markers, $\eta_\mathrm{pred}$ dashed,
    low-$\gamma^2$ bins faint}
  \caption{Measured strain-transfer and midpoint FRFs for [selected
    configurations]. [Describe: plateau, dip, amplification, phase flip.]}
  \label{fig:res:frf}
\end{figure}

\begin{table}[htb]
  \centering
  \caption{Resonance characterisation: measured $f_1$, quality factor $Q$,
    damping $\zeta$, and theoretical clamped-beam $f_1^\mathrm{th}$ per
    configuration. [From freqdomain resonance summary.]}
  \label{tab:res:resonance}
  % → build from freqdomain_eta.npz (f1_meas) + resonance_summary table
\end{table}

% ─────────────────────────────────────────────────────────────────────────────
\subsection{Method Validation and Estimator Robustness}
\label{sec:res:validation}
% Source: notebooks 02 + 03 §4–5 → timedomain_example.npz, freqdomain_frf.npz
% Short section; alternatively distribute these items into the sections above
% or move to an appendix — but one compact "the numbers are trustworthy"
% section reads well before the more exploratory sections below.

\begin{itemize}
  \item Example time-domain window (one cable, one frequency): reference
    elongation vs.\ cable elongation from all three strain methods, and
    $\eta(t)$ — the "how a single $\eta$ value is made" illustration.
  \item Agreement of the three strain estimators (gradient / arc-length /
    Fourier) across frequencies — quantify (e.g.\ relative spread) rather
    than plot for all cases.
  \item FRF estimator comparison ($H_1$/$H_2$/$H_v$/direct + coherence) for
    one dataset: estimators coincide in coherent bands, diverge exactly where
    $\gamma^2$ drops → the band-mean $\eta$ is estimator-insensitive
    (report the typical $H_1$-vs-$H_2$ difference within the quasi-static
    band).
  \item STFT single-frequency cross-check points overlaid on the Welch
    curves: independent windowing route, same FRF.
  \item Welch vs.\ time-domain medians (linearity segments, notebook 06 §5
    table) as a final consistency check between the two estimation routes.
\end{itemize}

\begin{figure}[htb]
  \centering
  \figplaceholder{timedomain\_example\_window (exists in Figures\_paper/)\\
    (a) $\Delta u_\mathrm{ends}$, $\delta x_\ell$ (M1/M2/M3);
    (b) $\eta(t)$ + median}
  \caption{Example analysis window. [Describe: method agreement, stable
    $\eta(t)$.]}
  \label{fig:res:window}
\end{figure}

\begin{figure}[htb]
  \centering
  \figplaceholder{estimator\_comparison (exists in Figures\_paper/)\\
    $|H_\eta|$ for $H_1$/$H_2$/$H_v$ + direct, coherence panel below}
  \caption{FRF estimator comparison. [Describe: bracket behaviour,
    insensitivity in the gated band.]}
  \label{fig:res:estimators}
\end{figure}

% ─────────────────────────────────────────────────────────────────────────────
\subsection{Induced Displacement and Strain Amplitudes}
\label{sec:res:amplitudes}
% Source: notebook 04 → amplitude_reference.npz

\begin{itemize}
  \item Band-mean ($<$25 Hz) endpoint displacement $\Delta u_\mathrm{ends}$
    and strain $\varepsilon_\mathrm{ends}$ per configuration — the reference
    values handed to the numerical simulations (table + overview figure).
  \item Amplitude vs.\ frequency overlays (1–500 Hz): where amplitudes peak
    (resonances), how the sag variants differ from taut ones.
  \item Shaker-to-cable contact: first-sensor displacement vs.\ imposed
    $\delta L$ (1:1 plot) and transfer $T(f)$ — where drive amplitude is
    lost before it even reaches the cable, and where local dynamic
    amplification occurs ($T>1$).
  \item Typical strain magnitudes in µε — relevant for arguing what a real
    DAS interrogator would resolve.
\end{itemize}

\begin{figure}[htb]
  \centering
  \figplaceholder{amplitude\_strain\_overview (exists in Figures\_paper/)\\
    band-mean $\varepsilon_\mathrm{ends}$ [µε] per configuration, log y}
  \caption{Quasi-static endpoint strain per configuration. [Describe: range,
    cable/gap trends.]}
  \label{fig:res:amplitudes}
\end{figure}

% Optional: shaker-transfer two-panel figure (u_first vs δL + T(f)) if the
% first-contact loss turns out to be a story worth telling.

% ─────────────────────────────────────────────────────────────────────────────
\subsection{Amplitude Linearity}
\label{sec:res:linearity}
% Source: notebook 06 → linearity_eta.npz

\begin{itemize}
  \item $\eta$ vs.\ instantaneous drive amplitude along the 0→1 ramps, per
    test frequency and cable: flat = linear coupling; slope/hysteresis =
    amplitude dependence (candidate mechanism: frictional slip at mounts).
  \item Report per cable/frequency: over which amplitude range $\eta$ is
    constant, and the threshold (if any) where behaviour changes.
  \item Envelope-ratio slip indicator: constant ratio along ramp = slip-free.
  \item Consequence for the main analysis: does the sweep amplitude sit in
    the linear regime (i.e.\ are the $\eta$–$\Theta$ results
    amplitude-independent)?
\end{itemize}

\begin{figure}[htb]
  \centering
  \figplaceholder{linearity: binned-median $\eta$ vs.\ $|\delta L|$,\\
    one panel per frequency or overlay per cable\\
    (from notebook 06, plot\_linearity\_eta\_vs\_amplitude)}
  \caption{Coupling efficiency against drive amplitude. [Describe: linearity
    range, thresholds.]}
  \label{fig:res:linearity}
\end{figure}

% ─────────────────────────────────────────────────────────────────────────────
\subsection{Extension--Compression Asymmetry}
\label{sec:res:asymmetry}
% Source: notebook 04 §6 → asymmetry.npz. Direct link to the numerical
% simulations that predict different behaviour for the two directions.

\begin{itemize}
  \item Motivation sentence tying to the simulations: sagged segment stiffens
    under extension, softens under compression.
  \item Half-cycle efficiencies $\eta_\mathrm{ext}$ vs.\
    $\eta_\mathrm{comp}$: band-mean scatter across all configurations
    (1:1 plot) — do sagged configurations depart from the diagonal while taut
    ones sit on it?
  \item Peak asymmetry $\alpha$ and even-harmonic ratio $R_2 = |Y(2f)|/|Y(f)|$
    vs.\ frequency for selected (taut vs.\ sagged) pairs; $R_3$ alongside for
    attribution (asymmetric/quadratic vs.\ symmetric/cubic nonlinearity).
  \item Amplitude-resolved asymmetry from the linearity ramps: does
    $\eta_\mathrm{ext}$–$\eta_\mathrm{comp}$ separation grow with drive
    amplitude?
  \item Honest negative-result framing if the effect is below the noise: an
    upper bound on the asymmetry is itself a useful statement for the
    simulation comparison.
\end{itemize}

\begin{figure}[htb]
  \centering
  \figplaceholder{asymmetry overview: $\eta_\mathrm{ext}$ vs.\
    $\eta_\mathrm{comp}$ band means,\\ 1:1 line, all configurations
    (from plot\_asymmetry\_overview / asymmetry.npz)}
  \caption{Extension vs.\ compression coupling efficiency. [Describe:
    departure from symmetry for sagged configurations, or upper bound.]}
  \label{fig:res:asymmetry}
\end{figure}

% ─────────────────────────────────────────────────────────────────────────────
\subsection{Modal Behaviour and Excitation Directionality}
\label{sec:res:modal}
% Source: notebook 05 → modal_summary.npz/.csv

\begin{itemize}
  \item Detected resonances across configurations (mode-frequency chart):
    frequencies, assigned mode orders, MAC values, polarization classes.
  \item $f_n$ vs.\ gap length against the Euler–Bernoulli $L^{-2}$ scaling
    and the $(\alpha_n/\alpha_1)^2$ frequency ratios — where the
    doubly-clamped beam model holds and where it does not.
  \item One deep-dive example: measured mode shape vs.\ best theoretical
    shape + MAC heatmap (shows the reader what a "good" match looks like).
  \item Shaker input bias: directionality spectra (x/y/z power fractions of
    the shaker head) next to detected peaks; raw vs.\ input-normalised
    detection — which "resonances" survive normalisation and which were
    excitation coloration. This resolves the puzzle of similar resonance
    frequencies across very different cables.
  \item Polarization statistics: fraction linear/elliptical/circular, and
    whether circular/elliptical modes correlate with off-axis input bands.
  \item Connect back to Section~\ref{sec:res:dynamic}: does the fingerprint
    $f_1$ agree with the FRF-based $f_1$?
\end{itemize}

\begin{figure}[htb]
  \centering
  \figplaceholder{shaker directionality: head spectra + power fractions\\
    + cable detection signal with peaks
    (from plot\_shaker\_directionality)}
  \caption{Shaker excitation directionality vs.\ detected cable resonances.
    [Describe: off-axis bands, coinciding peaks.]}
  \label{fig:res:shakerbias}
\end{figure}

\begin{figure}[htb]
  \centering
  \figplaceholder{mode-frequency chart across datasets\\ +
    example mode shape with theory overlay and MAC
    (from modal\_plotting)}
  \caption{Detected resonances and mode identification across
    configurations. [Describe: mode orders, MAC quality, gap scaling.]}
  \label{fig:res:modes}
\end{figure}

% ─────────────────────────────────────────────────────────────────────────────
% Closing paragraph idea (leads into Discussion): one short synthesis —
% geometry (Θ) predicts the quasi-static transfer within [x] %, the dynamic
% behaviour is organised by a single resonance whose identity/excitation is
% clarified by the modal + directionality analysis, and the model assumptions
% (linearity, symmetry) hold within [bounds] over the tested amplitude range.
